\begin{tikzpicture}[sb, x=0.18cm, y=0.62cm,
    rblk/.style={draw=sbInk!80, line width=0.5pt, fill=sbAux!55},
    sblk/.style={draw=sbInk!80, line width=0.5pt, fill=sbAux!20},
    lab/.style={font=\footnotesize, inner sep=1pt}]
  \def\p{67}\def\Q{4}
  \pgfplotstableread{\sbdata{jdegrees.dat}}\sbJTab
  \pgfplotstablegetrowsof{\sbJTab}%
  \pgfmathtruncatemacro\sbJLast{\pgfplotsretval-1}%
  \foreach \r in {0,...,\Q}{%
    \draw[sbOther!60, line width=0.3pt] (0,\r) rectangle (\p,\r+1);
    \node[lab, anchor=east] at (-0.6,\r+0.5) {$\r$};}
  \foreach \c in {8,16,...,64}
    \draw[sbOther!35, line width=0.3pt] (\c,0) -- (\c,\Q+1);
  \fill[sbOther!18] (64,0) rectangle (\p,\Q+1);
  \foreach \i in {0,...,\sbJLast}{%
    \pgfplotstablegetelem{\i}{kind}\of\sbJTab \let\kk\pgfplotsretval
    \pgfplotstablegetelem{\i}{j}\of\sbJTab \let\jj\pgfplotsretval
    \pgfplotstablegetelem{\i}{quotient}\of\sbJTab \let\qq\pgfplotsretval
    \pgfplotstablegetelem{\i}{reslo}\of\sbJTab \let\lo\pgfplotsretval
    \pgfplotstablegetelem{\i}{reshi}\of\sbJTab \let\hi\pgfplotsretval
    \pgfmathtruncatemacro\hip{\hi+1}%
    \if r\kk
      \draw[rblk] (\lo,\qq+0.12) rectangle (\hip,\qq+0.88);
      \ifnum\jj=0 \node[lab] at ({(\lo+\hip)/2},\qq+0.5) {$r_0$};
      \else \node[lab] at ({(\lo+\hip)/2},\qq+0.5) {\ifnum\jj=1 $r_1X^{p}$\else$r_{\jj}X^{\jj p}$\fi};\fi
    \else
      \draw[sblk] (\lo,\qq+0.12) rectangle (\hip,\qq+0.88);
      \fill[pattern={north east lines}, pattern color=sbOther]
        (\lo,\qq+0.12) rectangle (\hip,\qq+0.88);
      \ifnum\jj=0 \node[lab, fill=white, inner sep=1.5pt]
          at ({(\lo+\hip)/2},\qq+0.5) {$s_0f^c$};
      \else \node[lab, fill=white, inner sep=1.5pt]
          at ({(\lo+\hip)/2},\qq+0.5) {\ifnum\jj=1 $s_1f^cX^{p}$\else$s_{\jj}f^cX^{\jj p}$\fi};\fi
    \fi}
  \node[lab, rotate=90, anchor=south] at (-3.2,{(\Q+1)/2})
    {quotient by $p$};
  \foreach \c in {0,8,...,64}{%
    \draw[sbInk] (\c,0) -- (\c,-0.12);
    \node[lab, below] at (\c,-0.12) {$\c$};}
  \draw[sbInk] (\p,0) -- (\p,-0.12) node[lab, below] {$\p$};
  \node[lab, anchor=north] at ({\p/2},-0.75) {residue of the degree modulo $p$};
  \draw[decorate, decoration={brace, amplitude=4pt}, sbInk]
    (0,\Q+1.1) -- (32,\Q+1.1)
    node[midway, above=5pt, lab, align=center]
    {$\deg r_j\le D-1=31$};
  \draw[decorate, decoration={brace, amplitude=4pt}, sbInk]
    (32,\Q+1.1) -- (64,\Q+1.1)
    node[midway, above=5pt, lab, align=center]
    {$\deg(s_jf^cX^{jp})=(j+1)p+D+\deg s_j$, as $cd=p+D$};
  \node[lab, anchor=west, align=left, text=sbInk!80] at (\p+0.8,{(\Q+1)/2})
    {residues\\$2D,\dots,p-1$\\never used:\\$2D-1=63<p$};
  \node[font=\footnotesize, anchor=north, align=center] at ({\p/2},-1.45) {%
    \tikz\draw[rblk](0,0) rectangle (0.4cm,0.22cm); possible degrees of
    $r_jX^{jp}$ (row $j$)\qquad
    \tikz{\draw[sblk](0,0) rectangle (0.4cm,0.22cm);
      \fill[pattern={north east lines}, pattern color=sbOther]
      (0,0) rectangle (0.4cm,0.22cm);} possible degrees of $s_jf^cX^{jp}$
    (row $j+1$)};
\end{tikzpicture}
\caption{Why the auxiliary polynomial $r$ of Definition~\ref{def:aux-poly}
is nonzero (Proposition~\ref{prop:J-nonvanishing}), for $p=67$, $d=3$, so
$D=32$, $c=33$, $J=4$. The degree of each nonzero term of
$\sum_j(r_j+s_jf^c)X^{jp}$ is placed by its quotient (row) and residue
(column) modulo $p$; it lies in the block shown. Terms $r_jX^{jp}$ (solid
sky blue) lie in row $j$, residues $0$ to $D-1$; terms $s_jf^cX^{jp}$
(hatched) lie in row $j+1$, residues $D$ to $2D-1<p$; residues $2D$ to
$p-1$ (grey) are never used. Two terms of the same kind differ in the row,
two of different kinds in the column, so distinct nonzero terms have
distinct degrees and cannot cancel. Compare Figure~\ref{fig:residues}:
degrees at infinity play the role of pole orders at $P_0$.}
\label{fig:J-degrees}
